\chapter{Results}
\label{ch:results}

This chapter presents the main empirical findings. Section~\ref{sec:sumstats} describes the data. Section~\ref{sec:main_results} reports the headline cohort wage gap estimates. Section~\ref{sec:first_stage} presents the first-stage college attainment effects and implied returns. Section~\ref{sec:age_placebo} validates the identifying assumption through the age-cohort placebo. Section~\ref{sec:heterogeneity} examines heterogeneity by gender and race. Section~\ref{sec:robustness_results} collects robustness checks. Section~\ref{sec:pooled_results} briefly discusses the original pooled estimates.

\section{Summary Statistics}
\label{sec:sumstats}

Table~\ref{tab:sumstats} presents descriptive statistics for the analysis sample at the microregion-year level. The sample covers [FILL] microregions over the period 2002--2019. The average cumulative ProUni dose ($D_{r,2019}$) is [FILL] scholarships per 1,000 youth, with substantial geographic variation (SD = [FILL]). Low-dose microregions (20th--50th percentile) have an average dose of [FILL], while high-dose microregions ($>$75th percentile) average [FILL].

Average log wages (in minimum-wage units) for the full 22--35 sample are [FILL] (SD = [FILL]). The young cohort (22--28) earns [FILL] on average, compared to [FILL] for the old cohort (29--35), reflecting standard experience premia. The college-educated share is [FILL]\% among young workers and [FILL]\% among old workers.

% \begin{table}[htbp]
% \centering
% \caption{Summary Statistics}
% \label{tab:sumstats}
% \input{Tables_Final/table_sumstats.tex}
% \end{table}

\section{Main Results: Cohort Wage Gap}
\label{sec:main_results}

Figure~\ref{fig:headline} presents the event-study estimates for the cohort wage gap (young minus old log wages) separately for low-dose and high-dose microregions. Three features stand out.

First, pre-treatment coefficients are close to zero and jointly insignificant (F-test $p$ = [FILL] for low dose, [FILL] for high dose), supporting the parallel trends assumption.

Second, the low-dose ATT is positive and statistically significant at [FILL] (SE = [FILL]), indicating that ProUni raised the relative wage of the exposed cohort by approximately [FILL] log points. This corresponds to a [FILL]\% wage premium.

Third, the high-dose ATT is [FILL] (SE = [FILL]), which is [FILL: smaller/larger/similar]. The contrast between low and high dose is consistent with a saturation or diminishing-returns pattern: where ProUni scholarships are abundant, the marginal graduate faces more competition from other degree holders, compressing the college wage premium.

\begin{figure}[htbp]
    \centering
    \includegraphics[width=0.9\textwidth]{Figures_Final/fig_cohort_gap.png}
    \caption{Event-Study: Cohort Wage Gap (Young 22--28 vs.\ Old 29--35), by Dose Bin}
    \label{fig:headline}
\end{figure}

Table~\ref{tab:main_att} reports the overall ATT estimates for all cohort-based models.

% \begin{table}[htbp]
% \centering
% \caption{ATT Estimates: Cohort-Based Models}
% \label{tab:main_att}
% \input{Tables_Final/table_master_att.tex}
% \end{table}

\section{First Stage and Implied Returns}
\label{sec:first_stage}

To recover the individual-level return to a ProUni degree, I estimate the ``first stage'' effect of ProUni intensity on the college attainment gap (difference in college-educated worker shares between young and old cohorts). Figure~\ref{fig:first_stage} shows the event-study for this outcome.

The first-stage ATT is [FILL] (SE = [FILL]) for low dose and [FILL] (SE = [FILL]) for high dose. This indicates that ProUni increased the relative college share of the young cohort by [FILL] percentage points in low-dose areas.

The implied Wald/LATE estimate of the return to a ProUni degree is:
\[
    \hat{\rho}_{\text{LATE}}^{\text{Low}} = \frac{[FILL]}{[FILL]} = [FILL]
\]
with a delta-method standard error of [FILL]. This estimate is [FILL: consistent with / higher than / lower than] the Mincerian returns to higher education in Brazil documented by \citeA{BarbosaFilhoPessoa2008} and \citeA{MenezesFilho2001}.

\begin{figure}[htbp]
    \centering
    \includegraphics[width=0.9\textwidth]{Figures_Final/fig_first_stage.png}
    \caption{Event-Study: First Stage --- College Share Gap (Young $-$ Old)}
    \label{fig:first_stage}
\end{figure}

\section{Age-Specific Placebo}
\label{sec:age_placebo}

Figure~\ref{fig:age_placebo} plots the event-study estimates separately for the young (22--28) and old (29--35) cohorts within the same microregion-year cells. Under the identifying assumption, ProUni should affect only the young cohort; the old-cohort ATT should be approximately zero.

The young-cohort ATT is [FILL] (SE = [FILL]), while the old-cohort ATT is [FILL] (SE = [FILL]). The old-cohort estimate is [FILL: statistically indistinguishable from zero / small], confirming that the treatment effect is age-specific and not driven by general local economic conditions.

\begin{figure}[htbp]
    \centering
    \includegraphics[width=0.9\textwidth]{Figures_Final/fig_cohort_young_vs_old.png}
    \caption{Age-Specific Effects: Young (Exposed) vs.\ Old (Placebo), Low Dose}
    \label{fig:age_placebo}
\end{figure}

\section{Heterogeneity}
\label{sec:heterogeneity}

\subsection{Gender}

Figure~\ref{fig:gender} presents the cohort wage gap estimates separately for male and female workers. [FILL: Describe pattern --- e.g., ``Female workers show a larger cohort wage premium than males, consistent with ProUni's disproportionate enrollment of women and the higher marginal return to education for women in the Brazilian labor market.'']

\begin{figure}[htbp]
    \centering
    \includegraphics[width=0.48\textwidth]{Figures_Final/fig_cohort_gender_low.png}
    \hfill
    \includegraphics[width=0.48\textwidth]{Figures_Final/fig_cohort_gender_high.png}
    \caption{Gender Heterogeneity in Cohort Wage Premium}
    \label{fig:gender}
\end{figure}

\subsection{Race}

Figure~\ref{fig:race} compares the cohort wage gap for non-white and white workers. [FILL: Describe pattern --- e.g., ``Non-white workers show a larger treatment effect, consistent with ProUni's affirmative action mandate and the higher marginal return to degree attainment for historically disadvantaged groups.'']

\begin{figure}[htbp]
    \centering
    \includegraphics[width=0.48\textwidth]{Figures_Final/fig_cohort_race_low.png}
    \hfill
    \includegraphics[width=0.48\textwidth]{Figures_Final/fig_cohort_race_high.png}
    \caption{Racial Heterogeneity in Cohort Wage Premium}
    \label{fig:race}
\end{figure}

\section{Robustness}
\label{sec:robustness_results}

\subsection{TWFE Benchmark}

Table~\ref{tab:twfe} reports the standard two-way fixed effects estimates for the cohort wage gap and young-only wages. The TWFE coefficients are [FILL: qualitatively similar to / smaller than / larger than] the CS estimates, [FILL: consistent with the known bias of TWFE under heterogeneous treatment effects].

% \begin{table}[htbp]
% \centering
% \caption{TWFE Benchmark}
% \label{tab:twfe}
% \input{Tables_Final/table_twfe_cohort.csv}
% \end{table}

\subsection{Alternative Cutoffs}

Using 25th/75th and 33rd/67th percentile cutoffs instead of 20th/50th/75th yields ATT estimates of [FILL] and [FILL], respectively. These are [FILL: quantitatively similar], confirming that the main findings are not driven by the choice of discretization thresholds.

\subsection{Buffer-Zone Exclusion}

Excluding microregions in the middle of the dose distribution (between Low and High bins) yields [FILL: similar / slightly stronger] estimates, suggesting that SUTVA violations from spatial spillovers are unlikely to drive the results.

\subsection{Placebo Cohort (Ages 36--50)}

An external placebo test using workers aged 36--50 --- who should be entirely unaffected by ProUni --- yields ATT estimates of [FILL], confirming that the treatment effect is confined to the age range plausibly exposed to the program.

\subsection{Continuous Treatment (contdid)}

[FILL: If contdid succeeded, describe the dose-response curve. If it failed, note that the discretized results are robust to alternative cutoffs and reference the contdid package.]

\section{Original Pooled Estimates}
\label{sec:pooled_results}

For completeness, the pooled estimates for all workers aged 22--35 (without the cohort decomposition) show an overall ATT of [FILL] for low dose and [FILL] for high dose. These are [FILL: smaller than the cohort wage gap estimates], which is expected: the pooled outcome mixes the treated (young) and untreated (old) cohorts, diluting the signal. The cohort decomposition reveals the underlying individual-level effect that is masked by this aggregation.
